\section{The inequalities of Hölder and Minkowski}

In this No.~and in the following one, X and P have the same meaning as in No.~1, and M denotes a function defined on P that satisfies the conditions listed in No.~1.

PROPOSITION 2. --- Let \(\alpha\) and \(\beta\) be two numbers such that \(0 < \alpha < 1\) , \(0 < \beta < 1\) , \(\alpha + \beta = 1\) . If \(f\) and \(g\) are two finite functions \(\geqslant 0\) defined on X, and if \(M(f)\) and \(M(g)\) are finite, then

\[
M (f ^ {\alpha} g ^ {\beta}) \leqslant (M (f)) ^ {\alpha} (M (g)) ^ {\beta}\tag{5}
\]

(Hölder's inequality).

By Prop. 1, it all comes down to proving that, in \(\mathbf{R}^2\) , the set defined by the relations \(t_1 \geqslant 0\) , \(t_2 \geqslant 0\) , \(t_1^\alpha t_2^\beta \geqslant 1\) is convex, or again (FRV, I, §4, No.~1, Def. 1) that the function \(u(t) = t^{-\alpha/\beta}\) is convex for \(0 < t < +\infty\) . Now, setting \(r = \alpha/\beta\) , we have \(\mathrm{D}^2 u(t) = r(r + 1)t^{-r - 2}\) and, since \(r > 0\) , \(\mathrm{D}^2 u(t) > 0\) on \(]0, +\infty[\) , which proves the proposition (FRV, I, §4, No.~4, Cor. of Prop. 8).

COROLLARY. --- Let \(\alpha_{i}\) ( \(1 \leqslant i \leqslant n\) ) be n numbers \(\geqslant 0\) such that \(\sum_{i=1}^{n} \alpha_{i} = 1\) , and let \(f_{i}\) ( \(1 \leqslant i \leqslant n\) ) be n functions \(\geqslant 0\) defined on X, such that \(M(f_{i})\) is finite for \(1 \leqslant i \leqslant n\) . Under these conditions,

\[
M \left(f _ {1} ^ {\alpha_ {1}} f _ {2} ^ {\alpha_ {2}} \dots f _ {n} ^ {\alpha_ {n}}\right) \leqslant \left(M (f _ {1})\right) ^ {\alpha_ {1}} \left(M (f _ {2})\right) ^ {\alpha_ {2}} \dots \left(M (f _ {n})\right) ^ {\alpha_ {n}}.\tag{6}
\]

We can restrict attention to the case that \(\alpha_{i} > 0\) for every i. It suffices to argue by induction on n, by applying the inequality (5) to the numbers \(\alpha = \alpha_{1}\) and \(\beta = \sum_{i=2}^{n} \alpha_{i}\) , and to the functions \(f = f_{1}\) , \(g = (f_{2}^{\alpha_{2}} f_{3}^{\alpha_{3}} \cdots f_{n}^{\alpha_{n}})^{1/\beta}\) .

PROPOSITION 3. --- Let p be a real number \(\geqslant 1\) . If f and g are two finite functions \(\geqslant 0\) defined on X, then

\[
\left(M \big ((f + g) ^ {p} \big)\right) ^ {1 / p} \leqslant \left(M (f ^ {p})\right) ^ {1 / p} + \left(M (g ^ {p})\right) ^ {1 / p}\tag{7}
\]

(Minkowski's inequality).

We can restrict attention to the case that \(M(f^{p})\) and \(M(g^{p})\) are finite. By Prop. 1, we are reduced to proving that the set in \(R^{2}\) defined by the relations \(t_{1} \geqslant 0\) , \(t_{2} \geqslant 0\) , \(t_{1}^{1/p} + t_{2}^{1/p} \geqslant 1\) is convex, or again that the function \(u(t) = (1 - t^{1/p})^{p}\) is convex for \(0 \leqslant t \leqslant 1\) . Now,

\[
\mathrm{D} ^ {2} u (t) = \left(1 - \frac {1}{p}\right) t ^ {1 / p - 2} (1 - t ^ {1 / p}) ^ {p - 2} \geqslant 0
\]

for \(0 < t \leqslant 1\) , whence the proposition.

